\documentclass[12pt]{article}

\usepackage[T1]{fontenc}
\usepackage[utf8]{inputenc}
\usepackage[spanish, es-nodecimaldot]{babel}
\usepackage{graphicx}
\usepackage[a4paper, margin=2.5cm]{geometry}
\usepackage{float}
\usepackage{caption}
\usepackage{placeins}
\usepackage[hidelinks]{hyperref}

\title{Reporte de Visualización de Gráficas CMAT}
\author{Proyecto SofiaGC}
\date{\today}

\newcommand{\annualfigure}[4]{%
  \subsection{#2 (#3)}
  \begin{figure}[H]
    \centering
    \includegraphics[width=0.9\textwidth]{#1}
    \caption{Relación entre asesorías y calificación en #2 durante #3.}
  \end{figure}
  \noindent\textbf{Lectura breve.} #4
  \par\medskip
}

\newcommand{\periodfigure}[3]{%
  \subsection{#2}
  \begin{figure}[H]
    \centering
    \includegraphics[width=0.9\textwidth]{#1}
    \caption{Promedio anual de calificación en #2 para el periodo 2020--2024, con comparación entre estudiantes con asesoría y sin asesoría.}
  \end{figure}
  \noindent\textbf{Lectura breve.} #3
  \par\medskip
}

\begin{document}

\maketitle
\tableofcontents
\newpage

\section{Introducción}
Este documento reúne las gráficas disponibles del proyecto CMAT en dos bloques complementarios. El primero presenta diagramas de dispersión por año, donde se observa la relación entre el número de asesorías y la calificación final de cada materia. El segundo bloque muestra gráficas de líneas con el promedio anual de calificación para el periodo 2020--2024, comparando estudiantes con asesoría y sin asesoría.

En las gráficas de dispersión aparece una línea de referencia correspondiente a la calificación mínima aprobatoria de 7.5 cuando esa marca está incluida en la imagen. La intención del reporte es descriptiva: identificar patrones visuales, concentración de puntos, dispersión y comportamiento general, sin introducir inferencias estadísticas que no estén sustentadas directamente por la visualización.

También conviene notar que no todas las materias aparecen en todos los años. El reporte respeta exactamente el conjunto de imágenes existentes en las carpetas del proyecto y conserva su organización temporal original.

\section{Gráficas por año}

\subsection*{Criterio de lectura}
En cada figura, el eje horizontal representa el número de asesorías y el eje vertical la calificación obtenida. Una concentración de puntos cerca de cero asesorías indica que buena parte de los registros corresponde a estudiantes que no asistieron o asistieron poco; la dispersión vertical permite ver si, aun con un mismo nivel de asesoría, hubo resultados académicos muy distintos.

\section{Año 2019}
\annualfigure{graficos CMAT/graficos CMAT/2019/Álgebra_Lineal_2019.png}{Álgebra Lineal}{2019}{La nube de puntos permite revisar cómo se distribuyen las calificaciones para distintos niveles de asesoría en el primer año del conjunto. Se aprecia una mezcla de resultados por encima y por debajo del umbral aprobatorio, con mayor densidad visual en los niveles bajos de asesoría.}
\annualfigure{graficos CMAT/graficos CMAT/2019/Cálculo_I_2019.png}{Cálculo I}{2019}{La gráfica muestra una alta concentración de observaciones en pocos niveles de asesoría y una dispersión vertical amplia en las calificaciones. Esto sugiere que, dentro del mismo año, coexistieron trayectorias académicas muy distintas entre estudiantes con participaciones similares en asesoría.}
\annualfigure{graficos CMAT/graficos CMAT/2019/Cálculo_II_2019.png}{Cálculo II}{2019}{En esta visualización se distinguen estudiantes con desempeños muy variados incluso cuando las asesorías son escasas. La presencia de puntos alrededor y arriba de la línea aprobatoria ayuda a identificar que el grupo no tuvo un comportamiento homogéneo.}
\annualfigure{graficos CMAT/graficos CMAT/2019/Ecuaciones_Diferenciales_Ordinarias_2019.png}{Ecuaciones Diferenciales Ordinarias}{2019}{La dispersión de este curso permite comparar si el aumento en asesorías coincide visualmente con concentraciones de notas más altas. Más que una relación lineal clara, la figura enfatiza la variedad de resultados y la importancia de observar casos extremos.}
\annualfigure{graficos CMAT/graficos CMAT/2019/Estadística_Inferencial_2019.png}{Estadística Inferencial}{2019}{La gráfica presenta un conjunto de puntos distribuido en distintos niveles de calificación, con especial interés en la cercanía al nivel de aprobación. Visualmente, conviene fijarse en si la mayor concentración ocurre con pocas asesorías y en qué rango de notas se agrupa.}
\annualfigure{graficos CMAT/graficos CMAT/2019/Estadística_Para_Ciencias_De_La_Salud_2019.png}{Estadística para Ciencias de la Salud}{2019}{El patrón general permite evaluar la heterogeneidad del grupo en términos de asesorías y resultado final. La figura sirve sobre todo para reconocer si hay concentración en pocas asesorías y si las notas altas aparecen distribuidas o aisladas.}
\annualfigure{graficos CMAT/graficos CMAT/2019/Estadística_Para_Ciencias_Sociales_2019.png}{Estadística para Ciencias Sociales}{2019}{Aquí se observa la relación entre asistencia a asesorías y calificación desde una perspectiva descriptiva. La nube de puntos deja ver dispersión suficiente para pensar en un grupo diverso, con estudiantes cercanos al límite aprobatorio y otros claramente por encima.}
\annualfigure{graficos CMAT/graficos CMAT/2019/Estadística_Para_Negocios_2019.png}{Estadística para Negocios}{2019}{La visualización ayuda a identificar cómo se reparten las notas a lo largo de distintos niveles de asesoría. Resulta útil para comparar si los puntos aprobatorios se concentran más en ciertas zonas o si la variación persiste en todo el rango observado.}
\annualfigure{graficos CMAT/graficos CMAT/2019/Matemáticas_Universitarias_2019.png}{Matemáticas Universitarias}{2019}{El comportamiento de la nube de puntos permite revisar la amplitud del desempeño académico durante 2019. Se trata de una gráfica útil para detectar concentración en asesorías bajas y la coexistencia de calificaciones altas y bajas en ese mismo segmento.}
\annualfigure{graficos CMAT/graficos CMAT/2019/Teoría_De_Matrices_2019.png}{Teoría de Matrices}{2019}{La figura resume de forma visual cómo se combinan asesorías y rendimiento en esta materia. Más que mostrar una tendencia única, destaca la distribución general de casos y la separación entre registros reprobatorios y aprobatorios.}
\FloatBarrier

\section{Año 2020}
\annualfigure{graficos CMAT/graficos CMAT/2020/Álgebra_Lineal_2020.png}{Álgebra Lineal}{2020}{La gráfica de 2020 permite contrastar si la concentración de estudiantes sigue ocurriendo en niveles bajos de asesoría. La dispersión de calificaciones sugiere que el rendimiento no puede explicarse solo por la cantidad de asesorías observada en la figura.}
\annualfigure{graficos CMAT/graficos CMAT/2020/Cálculo_I_2020.png}{Cálculo I}{2020}{La nube de puntos conserva un enfoque descriptivo claro: comparar densidades, extremos y cercanía a la línea aprobatoria. Visualmente resulta importante notar qué tan separados aparecen los casos con notas altas respecto de los casos con bajo desempeño.}
\annualfigure{graficos CMAT/graficos CMAT/2020/Cálculo_II_2020.png}{Cálculo II}{2020}{Este gráfico permite identificar si las asesorías más frecuentes coinciden con una franja de calificaciones más estable. Aun así, la variación interna del conjunto muestra que el grupo presenta comportamientos académicos diferenciados.}
\annualfigure{graficos CMAT/graficos CMAT/2020/Estadística_Inferencial_2020.png}{Estadística Inferencial}{2020}{La distribución visual de los puntos ayuda a reconocer cuántos registros se encuentran claramente arriba del umbral de aprobación. También deja ver si existe acumulación de observaciones en pocos valores de asesoría, rasgo común en este tipo de seguimiento.}
\annualfigure{graficos CMAT/graficos CMAT/2020/Estadística_Para_Ciencias_De_La_Salud_2020.png}{Estadística para Ciencias de la Salud}{2020}{La gráfica muestra la relación entre participación en asesorías y resultado final sin forzar una conclusión causal. Su principal valor está en exhibir la dispersión del grupo y en hacer visible la cercanía de varios casos al criterio de aprobación.}
\annualfigure{graficos CMAT/graficos CMAT/2020/Estadística_Para_Ciencias_Sociales_2020.png}{Estadística para Ciencias Sociales}{2020}{La figura permite una lectura rápida de la variabilidad de calificaciones en 2020. Se puede observar si predominan estudiantes con pocas asesorías y si dentro de ese subconjunto hay tanto resultados sobresalientes como resultados bajos.}
\annualfigure{graficos CMAT/graficos CMAT/2020/Estadística_Para_Negocios_2020.png}{Estadística para Negocios}{2020}{En esta materia, la nube de puntos facilita comparar zonas de mayor densidad con el comportamiento respecto a la línea aprobatoria. Es una imagen útil para detectar dispersiones amplias y posibles casos atípicos en el extremo alto o bajo de calificación.}
\annualfigure{graficos CMAT/graficos CMAT/2020/Matemáticas_Universitarias_2020.png}{Matemáticas Universitarias}{2020}{La visualización resume cómo se distribuyeron las observaciones de 2020 en torno a asesorías y notas finales. La lectura sugiere prestar atención a la acumulación de puntos en pocos niveles de asesoría y al ancho del rango de calificaciones.}
\annualfigure{graficos CMAT/graficos CMAT/2020/Teoría_De_Matrices_2020.png}{Teoría de Matrices}{2020}{Esta gráfica aporta un panorama general del desempeño académico en la materia durante el año. La presencia de distintos niveles de calificación dentro de rangos similares de asesoría refuerza la idea de una cohorte heterogénea.}
\FloatBarrier

\section{Año 2021}
\annualfigure{graficos CMAT/graficos CMAT/2021/Álgebra_Lineal_2021.png}{Álgebra Lineal}{2021}{La figura de 2021 permite observar si el patrón de concentración en asesorías bajas se mantiene respecto a años previos. La dispersión vertical sigue siendo clave para interpretar la diversidad de resultados entre estudiantes con niveles parecidos de acompañamiento.}
\annualfigure{graficos CMAT/graficos CMAT/2021/Cálculo_I_2021.png}{Cálculo I}{2021}{La visualización presenta una nube de puntos útil para comparar el rendimiento en torno a la línea de aprobación. Más que una trayectoria única, se aprecia un conjunto de casos que se distribuye en varios rangos de nota aun con pocos cambios en asesorías.}
\annualfigure{graficos CMAT/graficos CMAT/2021/Cálculo_II_2021.png}{Cálculo II}{2021}{En este gráfico se aprecia la relación entre número de asesorías y resultado final desde una perspectiva puramente descriptiva. La diversidad de puntos por encima y por debajo del umbral resalta la variabilidad del grupo.}
\annualfigure{graficos CMAT/graficos CMAT/2021/Estadística_Inferencial_2021.png}{Estadística Inferencial}{2021}{La figura sirve para ubicar la parte del grupo que logra calificaciones aprobatorias y la que permanece cerca o por debajo del límite. Además, ayuda a distinguir si las observaciones se concentran en pocos valores del eje horizontal.}
\annualfigure{graficos CMAT/graficos CMAT/2021/Estadística_Para_Ciencias_De_La_Salud_2021.png}{Estadística para Ciencias de la Salud}{2021}{La nube de puntos ofrece una lectura clara de la amplitud del rendimiento en la materia. Visualmente, interesa detectar si las calificaciones altas aparecen repartidas en varios niveles de asesoría o si se concentran en pocos casos.}
\annualfigure{graficos CMAT/graficos CMAT/2021/Estadística_Para_Ciencias_Sociales_2021.png}{Estadística para Ciencias Sociales}{2021}{Esta gráfica permite valorar la dispersión general del grupo y la cercanía de varios puntos al criterio de aprobación. También facilita revisar si los estudiantes con pocas asesorías forman la mayor parte de la muestra visible.}
\annualfigure{graficos CMAT/graficos CMAT/2021/Estadística_Para_Negocios_2021.png}{Estadística para Negocios}{2021}{La visualización muestra un comportamiento que puede compararse con otros años de la misma materia. Su lectura se centra en la distribución de notas, la densidad por nivel de asesoría y la identificación de resultados extremos.}
\annualfigure{graficos CMAT/graficos CMAT/2021/Matemáticas_Universitarias_2021.png}{Matemáticas Universitarias}{2021}{El gráfico resume de manera compacta la relación entre asesorías y calificación durante 2021. La coexistencia de notas distintas dentro de tramos semejantes de asesoría apunta a un patrón de desempeño diverso.}
\annualfigure{graficos CMAT/graficos CMAT/2021/Teoría_De_Matrices_2021.png}{Teoría de Matrices}{2021}{La nube de puntos ayuda a revisar de forma visual la estructura del grupo en la materia. Es especialmente útil para detectar concentración de casos, amplitud de la dispersión y posibles puntos alejados del comportamiento principal.}
\FloatBarrier

\section{Año 2022}
\annualfigure{graficos CMAT/graficos CMAT/2022/Álgebra_Lineal_2022.png}{Álgebra Lineal}{2022}{La gráfica presenta un panorama de cómo se combinaron asesorías y calificaciones en 2022. La lectura visual se beneficia de observar tanto la densidad cercana a pocas asesorías como la posición de los registros respecto al umbral aprobatorio.}
\annualfigure{graficos CMAT/graficos CMAT/2022/Cálculo_I_2022.png}{Cálculo I}{2022}{En esta materia, la figura permite revisar si existen zonas donde se concentran las calificaciones más altas. También deja ver si la dispersión continúa siendo amplia en los niveles bajos de asesoría, lo que hablaría de un grupo no homogéneo.}
\annualfigure{graficos CMAT/graficos CMAT/2022/Cálculo_II_2022.png}{Cálculo II}{2022}{La visualización ofrece una comparación clara entre número de asesorías y resultado final sin sugerir relaciones deterministas. Su principal aporte es mostrar la variedad de casos y la distancia visual entre aprobados y no aprobados.}
\annualfigure{graficos CMAT/graficos CMAT/2022/Estadística_Inferencial_2022.png}{Estadística Inferencial}{2022}{La nube de puntos permite identificar el rango general de calificaciones observado en la materia durante 2022. Resulta útil para distinguir si la mayor parte del grupo se ubica cerca del umbral o si predomina un desempeño más alejado de esa referencia.}
\annualfigure{graficos CMAT/graficos CMAT/2022/Estadística_Para_Ciencias_De_La_Salud_2022.png}{Estadística para Ciencias de la Salud}{2022}{La figura resume visualmente el comportamiento del grupo con base en dos variables simples. La dispersión observada ayuda a reconocer que estudiantes con niveles similares de asesoría pueden terminar con resultados académicos distintos.}
\annualfigure{graficos CMAT/graficos CMAT/2022/Estadística_Para_Ciencias_Sociales_2022.png}{Estadística para Ciencias Sociales}{2022}{Esta gráfica facilita detectar concentraciones y vacíos en la distribución de puntos. La lectura breve debe enfocarse en la amplitud del rango de notas y en la posición relativa de los registros frente a la línea de aprobación.}
\annualfigure{graficos CMAT/graficos CMAT/2022/Estadística_Para_Negocios_2022.png}{Estadística para Negocios}{2022}{El gráfico permite comparar la consistencia del rendimiento dentro del año. Si varios puntos se agrupan en pocos niveles de asesoría, la principal diferencia entre estudiantes se vuelve más visible a través del eje de calificación.}
\annualfigure{graficos CMAT/graficos CMAT/2022/Matemáticas_Universitarias_2022.png}{Matemáticas Universitarias}{2022}{La visualización de 2022 aporta un resumen compacto del comportamiento del grupo en la materia. Lo más relevante es reconocer la dispersión vertical y la eventual presencia de casos altos o bajos poco frecuentes.}
\annualfigure{graficos CMAT/graficos CMAT/2022/Teoría_De_Matrices_2022.png}{Teoría de Matrices}{2022}{La figura muestra la distribución general del desempeño académico respecto a las asesorías registradas. La relación visual no se presenta como lineal, sino como un conjunto de observaciones con distintos niveles de rendimiento.}
\FloatBarrier

\section{Año 2023}
\annualfigure{graficos CMAT/graficos CMAT/2023/Álgebra_Lineal_2023.png}{Álgebra Lineal}{2023}{La gráfica de 2023 mantiene el enfoque de observar densidad, dispersión y cercanía al umbral de aprobación. Esto permite comparar visualmente el comportamiento del grupo con años anteriores de la misma asignatura.}
\annualfigure{graficos CMAT/graficos CMAT/2023/Cálculo_I_2023.png}{Cálculo I}{2023}{La nube de puntos resume de forma clara cómo se distribuyeron las calificaciones para distintos niveles de asesoría. La lectura visual se centra en la amplitud del rango de notas y en la concentración de observaciones en el eje horizontal.}
\annualfigure{graficos CMAT/graficos CMAT/2023/Cálculo_II_2023.png}{Cálculo II}{2023}{En este caso, la figura ayuda a detectar si los puntos aprobatorios se agrupan en determinadas zonas o si aparecen a lo largo de varios niveles de asesoría. La dispersión también permite reconocer la diversidad interna del grupo.}
\annualfigure{graficos CMAT/graficos CMAT/2023/Estadística_Inferencial_2023.png}{Estadística Inferencial}{2023}{La visualización ofrece una lectura directa del rango de desempeño alcanzado en la materia. Su utilidad principal está en mostrar la distribución general de notas y la posible concentración de registros cerca de pocas asesorías.}
\annualfigure{graficos CMAT/graficos CMAT/2023/Estadística_Para_Ciencias_Sociales_2023.png}{Estadística para Ciencias Sociales}{2023}{La gráfica permite comparar visualmente la composición del grupo durante 2023. Es particularmente útil para revisar cuántos casos quedan próximos a la línea aprobatoria y qué tan dispersas aparecen las calificaciones en el resto del plano.}
\annualfigure{graficos CMAT/graficos CMAT/2023/Estadística_Para_Negocios_2023.png}{Estadística para Negocios}{2023}{Aquí se aprecia un panorama de la relación entre asesorías y calificación que prioriza la observación de patrones generales. La figura ayuda a identificar concentración de casos, dispersión y posibles valores atípicos.}
\annualfigure{graficos CMAT/graficos CMAT/2023/Matemáticas_Universitarias_2023.png}{Matemáticas Universitarias}{2023}{La nube de puntos permite revisar si el grupo presenta una distribución compacta o más bien extendida en términos de calificación. En ambos casos, la lectura destaca la heterogeneidad del rendimiento observado.}
\annualfigure{graficos CMAT/graficos CMAT/2023/Teoría_De_Matrices_2023.png}{Teoría de Matrices}{2023}{La figura resume la estructura visual del curso en 2023 y facilita comparaciones temporales posteriores. El peso descriptivo recae en la ubicación de los puntos, su concentración y la separación entre distintos rangos de nota.}
\FloatBarrier

\section{Año 2024}
\annualfigure{graficos CMAT/graficos CMAT/2024/Álgebra_Lineal_2024.png}{Álgebra Lineal}{2024}{La gráfica permite identificar rápidamente la parte del grupo que supera la línea aprobatoria y la que permanece debajo. También ofrece un panorama claro de cómo se distribuyen los estudiantes en los niveles bajos y medios de asesoría.}
\annualfigure{graficos CMAT/graficos CMAT/2024/Cálculo_I_2024.png}{Cálculo I}{2024}{En esta visualización se aprecia una fuerte concentración de casos en pocas o nulas asesorías, junto con una dispersión amplia de calificaciones. Esa combinación sugiere que, aun dentro del segmento mayoritario, el desempeño académico fue muy desigual.}
\annualfigure{graficos CMAT/graficos CMAT/2024/Cálculo_II_2024.png}{Cálculo II}{2024}{La figura de 2024 ayuda a examinar la relación entre acompañamiento académico y resultado final de manera comparativa. Más que confirmar una tendencia única, deja ver una variedad de trayectorias de desempeño dentro del curso.}
\annualfigure{graficos CMAT/graficos CMAT/2024/Estadística_Para_Ciencias_Sociales_2024.png}{Estadística para Ciencias Sociales}{2024}{La nube de puntos permite ubicar visualmente el peso relativo de los casos aprobatorios y no aprobatorios. La dispersión observada sugiere que el grupo no se concentra en un solo patrón de rendimiento.}
\annualfigure{graficos CMAT/graficos CMAT/2024/Estadística_Para_Negocios_2024.png}{Estadística para Negocios}{2024}{La gráfica muestra cómo se reparte el desempeño final a lo largo de distintos niveles de asesoría. Resulta útil para reconocer si la mayoría de los registros se agrupa en pocos valores del eje horizontal y qué tan amplio es el rango de notas.}
\annualfigure{graficos CMAT/graficos CMAT/2024/Matemáticas_Universitarias_2024.png}{Matemáticas Universitarias}{2024}{La visualización permite revisar de forma sintética la dispersión del grupo durante 2024. Las diferencias entre estudiantes se aprecian sobre todo en el eje de calificación, incluso cuando comparten niveles parecidos de asesoría.}
\FloatBarrier

\section{Año 2025}
\annualfigure{graficos CMAT/graficos CMAT/2025/Álgebra_Lineal_2025.png}{Álgebra Lineal}{2025}{La gráfica más reciente de esta materia permite una lectura preliminar del comportamiento del grupo. La distribución de puntos sigue siendo útil para distinguir concentración, dispersión y cercanía al umbral aprobatorio.}
\annualfigure{graficos CMAT/graficos CMAT/2025/Cálculo_I_2025.png}{Cálculo I}{2025}{La figura resume la relación entre asesorías y calificación en el año más reciente disponible. La interpretación visual debe centrarse en la acumulación de casos con pocas asesorías y en la amplitud del rango de resultados finales.}
\annualfigure{graficos CMAT/graficos CMAT/2025/Cálculo_II_2025.png}{Cálculo II}{2025}{Esta visualización permite detectar de manera inmediata si predominan los registros aprobatorios o si el grupo se encuentra más repartido alrededor del umbral. La variedad de puntos refleja que todavía hay comportamientos académicos diferenciados.}
\annualfigure{graficos CMAT/graficos CMAT/2025/Estadística_Para_Ciencias_Sociales_2025.png}{Estadística para Ciencias Sociales}{2025}{La nube de puntos presenta un panorama compacto del desempeño de la materia en 2025. Su lectura es útil para ubicar la dispersión de notas y la posible concentración de registros en pocas asesorías.}
\annualfigure{graficos CMAT/graficos CMAT/2025/Estadística_Para_Negocios_2025.png}{Estadística para Negocios}{2025}{La gráfica ofrece una referencia reciente para comparar con años anteriores de la asignatura. Visualmente, destaca la distribución de los casos y la separación entre resultados que quedan por encima o por debajo del criterio de aprobación.}
\annualfigure{graficos CMAT/graficos CMAT/2025/Matemáticas_Universitarias_2025.png}{Matemáticas Universitarias}{2025}{La figura cierra la serie anual mostrando el comportamiento más reciente del grupo. Como en otros cursos, la atención se centra en la densidad de observaciones y en la dispersión vertical de las calificaciones.}
\FloatBarrier

\section{Evolución en el periodo 2020--2024}

\subsection*{Criterio de lectura}
Las siguientes gráficas comparan dos series de promedio anual: estudiantes con asesoría y estudiantes sin asesoría. La lectura principal consiste en observar la separación entre ambas líneas, su estabilidad a través del tiempo y los años donde la diferencia parece ampliarse o reducirse.

\periodfigure{graficos de lineas/graficos de lineas/Promedio_Álgebra_Lineal_2020_2024.png}{Álgebra Lineal}{La comparación temporal permite identificar si existe una brecha consistente entre el promedio con asesoría y el promedio sin asesoría. Más allá de la magnitud exacta, la figura resulta útil para observar estabilidad, cambios graduales y posibles años de mayor acercamiento o separación.}
\periodfigure{graficos de lineas/graficos de lineas/Promedio_Cálculo_I_2020_2024.png}{Cálculo I}{La serie anual permite revisar la trayectoria del promedio en ambos grupos a lo largo del periodo 2020--2024. La lectura visual se concentra en el orden relativo de las líneas y en la variación de la distancia entre ellas de un año a otro.}
\periodfigure{graficos de lineas/graficos de lineas/Promedio_Cálculo_II_2020_2024.png}{Cálculo II}{La gráfica resume de forma clara la evolución del desempeño promedio en la materia durante cinco años. Su principal interés está en comparar la continuidad de las tendencias y verificar si la diferencia entre grupos se mantiene relativamente estable.}
\periodfigure{graficos de lineas/graficos de lineas/Promedio_Estadística_Inferencial_2020_2024.png}{Estadística Inferencial}{La visualización permite contrastar dos trayectorias promedio en el tiempo sin perder de vista el nivel general de calificaciones. Resulta útil para detectar años de mayor variación y para evaluar si las líneas convergen o divergen en ciertos momentos del periodo.}
\periodfigure{graficos de lineas/graficos de lineas/Promedio_Estadística_Para_Ciencias_De_La_Salud_2020_2024.png}{Estadística para Ciencias de la Salud}{La figura destaca la comparación entre promedios anuales con y sin asesoría, mostrando si ambos recorridos siguen patrones parecidos. La atención debe ponerse en la continuidad de las curvas y en la brecha observable entre los dos grupos.}
\periodfigure{graficos de lineas/graficos de lineas/Promedio_Estadística_Para_Ciencias_Sociales_2020_2024.png}{Estadística para Ciencias Sociales}{La lectura temporal de esta materia ayuda a reconocer si el comportamiento promedio cambia poco o mucho entre años. También permite revisar si la diferencia entre estudiantes con asesoría y sin asesoría permanece constante o fluctúa.}
\periodfigure{graficos de lineas/graficos de lineas/Promedio_Estadística_Para_Negocios_2020_2024.png}{Estadística para Negocios}{La gráfica facilita una comparación directa de trayectorias promedio en ambos grupos durante el periodo completo. Su valor está en mostrar de manera sintética continuidad, variaciones interanuales y distancia relativa entre líneas.}
\periodfigure{graficos de lineas/graficos de lineas/Promedio_Matemáticas_Universitarias_2020_2024.png}{Matemáticas Universitarias}{La visualización resume la evolución del promedio anual con una lectura sencilla y comparativa. La estabilidad o cambio de la brecha entre líneas constituye el elemento central para interpretar el comportamiento de la materia en el tiempo.}
\periodfigure{graficos de lineas/graficos de lineas/Promedio_Teoría_De_Matrices_2020_2024.png}{Teoría de Matrices}{La figura permite cerrar el bloque temporal observando si ambos promedios siguen recorridos similares o si presentan diferencias más marcadas en algunos años. El análisis visual debe priorizar la forma general de las líneas y la consistencia de la separación entre grupos.}
\FloatBarrier

\section{Conclusiones generales}
En conjunto, las gráficas por año muestran que la lectura más útil proviene de la densidad de observaciones, la dispersión vertical de las calificaciones y la cercanía de los casos al umbral aprobatorio. En varias materias se aprecia que una parte importante de los registros se concentra en niveles bajos de asesoría, mientras que el desempeño final sigue presentando una amplitud considerable.

Por su parte, las gráficas del periodo 2020--2024 complementan esa lectura con una visión agregada del comportamiento promedio. Estas figuras permiten comparar de forma directa los recorridos anuales de estudiantes con asesoría y sin asesoría, identificar estabilidad o cambios en la brecha entre grupos y ofrecer un panorama más compacto para contrastes entre materias.

En términos documentales, este reporte deja reunidas todas las visualizaciones disponibles en las carpetas activas del proyecto, con una explicación breve para cada una. Esto facilita tanto la consulta rápida por año como la revisión transversal de tendencias en el tiempo.

\end{document}
